\section{Related Work}

Our work connects three previously independent research areas: contamination detection, data attribution, and training data curation. We review each and highlight how their separation has obscured the curation-contamination relationship we address.

\subsection{Benchmark Contamination Detection}

Detecting test data contamination in LLM training corpora has become essential for reliable evaluation. \textbf{N-gram overlap methods} identify exact or near-exact matches between training and test examples. Recent work recommends $n=8$ as the detection threshold, noting that larger $n$ leads to false negatives while smaller $n$ generates excessive false positives \citep{contam2024}. The LLM-Decontaminator extends detection to rephrased samples through embedding similarity, addressing evasion via paraphrasing \citep{shi2023rethinking}.

\textbf{Membership inference} approaches operate without explicit training corpus access. Min-K\%++ identifies training samples as local probability maxima in model output distributions \citep{zhang2024minkpp}. CDD (Contamination Detection via output Distribution) complements detection with TED (Trustworthy Evaluation via output Distribution), achieving up to 66.9\% contamination mitigation \citep{dong2024cdd}. The Contamination Taxonomy \citep{palavalli2024taxonomy} categorizes contamination by type (input-only, input-output) and impact mechanism.

\textbf{Limitation:} These methods detect contamination presence/absence but do not explain how curation decisions affect contamination rates. Detection operates \textit{after} training; our work addresses the curation stage \textit{before} training where contamination is amplified or attenuated.

\subsection{Data Attribution Methods}

Data attribution identifies which training examples contribute to model predictions. \textbf{Influence functions} \citep{koh2017influence} approximate leave-one-out effects but scale poorly to large models. \textbf{TRAK} (Attributing Model Behavior at Scale) \citep{park2023trak} enables practical attribution through random projection, achieving 6500$\times$ speedup via the LoGra algorithm \citep{choe2024logra}. DDA addresses fitting error considerations in attribution estimation \citep{wu2024dda}.

Attribution has been applied to data valuation, detecting mislabeled examples, and understanding model behavior. However, it has not been combined with contamination detection to trace \textit{which} influential examples derive from benchmark overlap.

\textbf{Limitation:} Prior attribution work treats all high-influence examples uniformly. We distinguish between legitimate high-influence examples (valuable training signal) and contaminated high-influence examples (benchmark memorization), using CCR to quantify the distinction.

\subsection{Training Data Curation}

Modern LLM training relies heavily on data curation. \textbf{DataComp-LM} \citep{li2024datacomp} established that model-based filtering significantly affects downstream performance, with DCLM-Baseline achieving strong results through careful selection. \textbf{SlimPajama} \citep{shen2023slimpajama} demonstrated that global versus local deduplication strategies yield different model capabilities. \textbf{FineWeb} pipelines \citep{penedo2025fineweb} combine multiple filtering stages (perplexity, classifier, heuristics) for web-scale curation.

Perplexity-based filtering---selecting documents that score well against reference language models---has emerged as a core technique. RedPajama-V2 includes pre-computed \texttt{ccnet\_perplexity} scores using Wikipedia-trained models, enabling researchers to filter by perplexity bucket.

\textbf{Limitation:} Curation research focuses on quality-performance relationships without analyzing contamination effects. DataComp-LM showed filtering improves benchmark scores but did not investigate whether improvements reflect genuine capability or contamination amplification. Our work fills this gap.

\subsection{Connecting the Three Areas}

The gap we address is visualized in Figure~\ref{fig:ccr_scaling}: contamination detection identifies problematic examples, attribution traces influence, curation shapes corpus composition---but no prior work connects all three. This separation means practitioners cannot answer: ``How much does my curation strategy amplify benchmark contamination, and does it matter for model capabilities?''

We introduce CCR as the connecting metric. By combining TRAK attribution with n-gram contamination detection, CCR quantifies the contamination-attributed fraction of benchmark performance. This enables direct comparison across curation strategies and causal validation through removal experiments.

Our approach differs from concurrent work on clean benchmarks (e.g., LiveCodeBench \citep{jain2024livecodebench}) which constructs new evaluation sets to avoid contamination. We instead \textit{measure} contamination effects to understand existing benchmarks, providing a complementary perspective on evaluation integrity.
